\chapter{Exotic classes}

\providecommand{\Frob}{\mathrm{Frob}}

This chapter is a faithful transcription of \S7 of Ancona's paper. We fix an
abelian variety \(A\) of dimension four over a finite field \(k\) (with a fixed
algebraic closure \(\bar k\)) and a CM-structure for \(A\)
(\cref{def:CM-structure}). By \cref{prop:cm-decomposition}(2) this produces
motives \(\mathcal{M}_I \in \CHM(k)_{\Q}\) which are direct factors of
\(\hh^{4}(A)\). Some of them, the \emph{exotic} ones, carry algebraic classes
but no Lefschetz classes; they are precisely the factors that the main theorem
\(\cref{thm:scht}\) cannot dispose of by the Lefschetz formalism and must instead
feed into the rank-two orthogonal-motive theorem. The aim of the chapter
(\cref{prop:exotic-dim2}) is to show that, after a finite extension of \(k\) and a
suitable choice of CM-structure, every exotic motive has dimension two.

Throughout we use \cref{def:cm-sigma}: the set \(\Sigma\), the rank-one factors
\(M_\sigma\) of \(\hh^1(A)\) from \cref{prop:decomposition}, and the action
\(\bar{\cdot}\) of complex conjugation on \(\Sigma\). Frobenius acts on each of
the eight \(M_\sigma\) (\cref{rem:frob}); we write \(\alpha_\sigma \in \bar\Q\)
for its eigenvalue on \(M_\sigma\), and use the same symbol \(\bar{\cdot}\) for
complex conjugation on \(\bar\Q\). Recall (\cref{prop:cm-decomposition}) that for
\(I \subset \Sigma\) the rank-one factor \(M_I = \bigotimes_{\sigma \in I}
M_\sigma\) carries the Frobenius eigenvalue \(\prod_{\sigma \in I} \alpha_\sigma\).

\section{Exotic classes}

\subsection*{Exotic motives and exotic classes}

\begin{definition}[Exotic motive]
  \label{def:exotic-motive}
  \lean{MR4199442StandardConjecturesForAbelianFourfolds.ExoticMotive}
  \uses{prop:cm-decomposition, def:lefschetz-class, def:num-cycles}
  % SOURCE: [Ancona], §7, Definition (definizione) (read from references/source/fourfolds_2020.tex)
  % SOURCE QUOTE: "Let  $\mathcal{M}_I \in \CHM (k)_{\Q}$ be a  direct factor of $\mathfrak{h}^{4}(A)$ as constructed in Proposition \ref{primadecomposizione}(\ref{decomposizioneQ}). The motive $\mathcal{M}_I $  is called exotic if the  space of algebraic classes $\Hom_{\NUM(k)_\Q}(\one, \mathcal{M}_I (2))$ is nonzero and it does not contain any  nonzero Lefschetz class (Definition \ref{defexotic}).
  % An element in the $\mathbb{Q}$-vector space $\Hom_{\NUM(k)_\Q}(\one, \mathcal{M}_I (2))$ will be called an exotic class."
  \textit{Source: Ancona, arXiv:1806.03216, §7.}
  Let \(\mathcal{M}_I \in \CHM(k)_{\Q}\) be a direct factor of \(\hh^{4}(A)\) as
  constructed in \cref{prop:cm-decomposition}(2). The motive \(\mathcal{M}_I\) is
  called \emph{exotic} if its space of algebraic classes
  \(\Hom_{\NUM(k)_\Q}(\one, \mathcal{M}_I(2))\) is nonzero and contains no nonzero
  Lefschetz class (\cref{def:lefschetz-class}). An element of the \(\Q\)-vector
  space \(\Hom_{\NUM(k)_\Q}(\one, \mathcal{M}_I(2))\) is called an \emph{exotic
  class}.
\end{definition}

\begin{remark}
  \label{rem:exotic-equal-dim}
  \uses{def:exotic-motive, prop:factor-algebraic}
  % SOURCE: [Ancona], §7, Remark (remarkuguale) (read from references/source/fourfolds_2020.tex)
  % SOURCE QUOTE: "Notice that, by Proposition \ref{talvez}(\ref{generatoalgebrico}), any exotic motive verifies the equality
  % \begin{equation} \label{eq:uguale}
  % \dim_\Q \Hom_{\NUM(k)_\Q}(\one, \mathcal{M}_I (2)) = \dim \mathcal{M}_I.
  %  \tag{$\#$}
  % \end{equation}
  % Examples of exotic classes on abelian fourfolds (and especially of those that cannot be lifted to characteristic zero) will be discussed in  Appendix \ref{geometricexamples}."
  \textit{Source: Ancona, arXiv:1806.03216, §7.}
  Since an exotic motive has a nonzero space of algebraic classes,
  \cref{prop:factor-algebraic}(1) applies to it and gives the equality
  \[\dim_\Q \Hom_{\NUM(k)_\Q}(\one, \mathcal{M}_I(2)) = \dim \mathcal{M}_I.\]
  Thus, for exotic motives, the dimension of the space of exotic classes equals
  the dimension of the motive.
  (The geometric examples of exotic classes on abelian fourfolds discussed in
  Appendix A of the reference are out of scope and are not formalized.)
\end{remark}

\begin{proposition}[Exotic motives have dimension two]
  \label{prop:exotic-dim2}
  \lean{MR4199442StandardConjecturesForAbelianFourfolds.exoticDimTwo}
  \uses{def:exotic-motive, lem:exotic-frobenius, def:exotic-subset,
    lem:exotic-partition, lem:exotic-count, lem:exotic-supersingular,
    lem:exotic-threefold-factor, lem:exotic-split, def:exotic-assumption}
  % SOURCE: [Ancona], §7, Proposition (exoticdim2) (read from references/source/fourfolds_2020.tex)
  % SOURCE QUOTE: "Suppose that $A$ has dimension four. Then there exist a finite extension $k'$ of $k$ and a  CM-structure for $A\times_k k'$ such that any exotic motive $\mathcal{M}_I \in\CHM (k')_{\Q}$  of $\mathfrak{h}^{4}(A\times_k k')$ has dimension two."
  \textit{Source: Ancona, arXiv:1806.03216, §7.}
  Suppose \(A\) has dimension four. Then there exist a finite extension \(k'\) of
  \(k\) and a CM-structure for \(A \times_k k'\) such that any exotic motive
  \(\mathcal{M}_I \in \CHM(k')_{\Q}\) of \(\hh^{4}(A \times_k k')\) has dimension
  two.
\end{proposition}

% SOURCE QUOTE PROOF: "This is a combination of the  last four lemmas."
\begin{proof}
  \uses{def:exotic-assumption, rem:exotic-assumption-lef, lem:exotic-count,
    lem:exotic-partition, lem:exotic-supersingular, lem:exotic-threefold-factor,
    lem:exotic-split, rem:exotic-equal-dim}
  The argument combines the four lemmas established in the remainder of this
  chapter; we sketch it here and prove the supporting lemmas below.
  After a finite extension of \(k\) we may assume \cref{def:exotic-assumption}
  holds (\cref{rem:exotic-assumption-lef}), and exotic motives stay exotic under
  further finite extension. By \cref{lem:exotic-count} there are zero, two or
  four exotic subsets, and by \cref{lem:exotic-partition}(3) together with the
  equality of \cref{rem:exotic-equal-dim} each exotic motive is even-dimensional.

  If there is no exotic subset there is nothing to prove. If there are two, they
  are a conjugate pair \(I, \bar I\), forming a single \(\Gal(\bar\Q/\Q)\)-orbit
  of cardinality two; hence the corresponding exotic motive \(\mathcal{M}_I\) has
  dimension two.

  If there are four, then by \cref{lem:exotic-supersingular}, after passing to
  \(\F_{q^2}\) the fourfold becomes isogenous to \(E \times X\) with \(E\) a
  supersingular elliptic curve on which Frobenius acts as \(q \cdot \id\); over
  \(\F_{q^2}\) such an \(E\) satisfies \(\dim_\Q \End(E) \otimes_\Z \Q = 4\). By
  \cref{lem:exotic-threefold-factor} the space of exotic classes lies in
  \(\mathcal{M}_{I \cap \Sigma_X} \otimes \hh^1(E)\) with
  \(\mathcal{M}_{I \cap \Sigma_X}\) of rank two and spanned by algebraic classes.
  Choosing the CM-structure on \(E\) provided by \cref{lem:exotic-split}, this
  motive splits as
  \(\mathcal{M}_{I \cap \Sigma_X} \otimes \hh^1(E) =
    \mathcal{M}_{(I \cap \Sigma_X) \cup \sigma} \oplus
    \mathcal{M}_{(I \cap \Sigma_X) \cup \bar\sigma}\),
  a sum of two motives of rank two. With this CM-structure each exotic motive is
  one of these rank-two summands (or a conjugate pair as above), hence has
  dimension two.

  Taking \(k'\) to be the finite extension realizing \cref{def:exotic-assumption}
  and the passage to \(\F_{q^2}\), and the CM-structure to be the one just
  selected, every exotic motive of \(\hh^{4}(A \times_k k')\) has dimension two.
\end{proof}

\subsection*{Frobenius relations and exotic subsets}

We now fix the standing hypothesis used in the lemmas below.

\begin{definition}[Algebraic classes defined over the base]
  \label{def:exotic-assumption}
  \lean{MR4199442StandardConjecturesForAbelianFourfolds.AlgebraicClassesOverBase}
  \uses{def:num-cycles}
  % SOURCE: [Ancona], §7, Assumption (assumassum) (read from references/source/fourfolds_2020.tex)
  % SOURCE QUOTE: "We will suppose that all the algebraic classes $\mathcal{Z}^*_\num (A_{\bar{k}})_\Q$ are defined over $k$."
  \textit{Source: Ancona, arXiv:1806.03216, §7.}
  We suppose that all the algebraic classes
  \(\mathcal{Z}^*_\num(A_{\bar k})_\Q\) are defined over \(k\).
\end{definition}

\begin{remark}
  \label{rem:exotic-assumption-lef}
  \uses{def:exotic-assumption, def:exotic-motive}
  % SOURCE: [Ancona], §7, Remark (assumptionlef) (read from references/source/fourfolds_2020.tex)
  % SOURCE QUOTE: "Note that the assumption above always holds after a finite extension of $k$. Notice also that, under this assumption, a class which becomes Lefschetz after a finite extension of the base field must be already Leftschetz. Hence, a motive which is exotic over the base field will still be exotic after a finite extension."
  \textit{Source: Ancona, arXiv:1806.03216, §7.}
  \cref{def:exotic-assumption} always holds after a finite extension of \(k\).
  Under it, a class that becomes Lefschetz after a finite extension of the base
  field is already Lefschetz; hence a motive that is exotic over \(k\) remains
  exotic after any finite extension.
\end{remark}

\begin{lemma}[Frobenius relations of an exotic motive]
  \label{lem:exotic-frobenius}
  \lean{MR4199442StandardConjecturesForAbelianFourfolds.exoticFrobeniusRelations}
  \uses{def:exotic-motive, prop:factor-algebraic, prop:cm-decomposition, rem:frob}
  % SOURCE: [Ancona], §7, Lemma (musica) (read from references/source/fourfolds_2020.tex)
  % SOURCE QUOTE: "Let $q$ be the cardinality of $k$ and let $\mathcal{M}_I$ be an exotic motive. Then we have  the relation
  % \begin{equation} \label{eq:1}
  % \prod_{\sigma \in I}\alpha_\sigma =q^2.
  % \end{equation}
  % Moreover, we have the property
  % \begin{equation} \label{eq:2}
  %  \alpha_\sigma \cdot \alpha_\tau \neq q, \, \, \forall \sigma \neq \tau, \,\, \sigma,\tau \in I.
  % \end{equation}"
  \textit{Source: Ancona, arXiv:1806.03216, §7.}
  Let \(q = \#k\) and let \(\mathcal{M}_I\) be an exotic motive. Then
  \[\prod_{\sigma \in I} \alpha_\sigma = q^2,\]
  and moreover
  \[\alpha_\sigma \cdot \alpha_\tau \neq q
    \qquad \text{for all distinct } \sigma, \tau \in I.\]
\end{lemma}

% SOURCE QUOTE PROOF: "By Proposition \ref{talvez}(\ref{generatoalgebrico}) the $\ell$-adic realization of $\mathcal{M}_I$ is spanned by algebraic cycles hence  Frobenius acts on it by multiplication by $q^2$.  On the other hand, by construction (see  Proposition \ref{primadecomposizione}(\ref{decomposizioneL})),  Frobenius acts on the line $ M_I\subset \mathcal{M}_I$ via the multiplication by $\prod_{\sigma \in I}\alpha_\sigma $. This gives (\ref{eq:1}).
% Suppose now that  (\ref{eq:2}) is not satisfied.    Then $(\ref{eq:1})$ would force the set    $\{\alpha_\sigma\}_{\sigma \in I}$ to be of the form  $\alpha,q/\alpha, \beta, q/\beta$.   Each of the pairs $(\alpha,q/\alpha)$ and $(\beta, q/\beta)$ would correspond to a Frobenius invariant class in $H^2_\ell (A) (1)$. As the Tate conjecture for divisors on abelian varieties is known \cite[Theorem 4]{Tate}, each of these pairs corresponds to the class of a divisor, hence   $M_I$ would contain a Lefschetz class."
\begin{proof}
  \uses{def:exotic-motive, prop:factor-algebraic, prop:cm-decomposition}
  By \cref{prop:factor-algebraic}(1) the \(\ell\)-adic realization of
  \(\mathcal{M}_I\) is spanned by algebraic cycles, so Frobenius acts on it by
  multiplication by \(q^2\). On the other hand, by the construction of the
  factors (\cref{prop:cm-decomposition}(1)), Frobenius acts on the line
  \(M_I \subset \mathcal{M}_I\) by multiplication by
  \(\prod_{\sigma \in I} \alpha_\sigma\). Comparing the two gives
  \(\prod_{\sigma \in I} \alpha_\sigma = q^2\).

  Suppose the second relation fails, say
  \(\alpha_\sigma \alpha_\tau = q\) for some distinct \(\sigma, \tau \in I\).
  Together with \(\prod_{\sigma \in I} \alpha_\sigma = q^2\) this forces the
  four eigenvalues \(\{\alpha_\sigma\}_{\sigma \in I}\) to be of the form
  \(\alpha, q/\alpha, \beta, q/\beta\). Each pair \((\alpha, q/\alpha)\) and
  \((\beta, q/\beta)\) cuts out a Frobenius-invariant class in
  \(H^2_\ell(A)(1)\). Since the Tate conjecture for divisors on abelian
  varieties is known \cite[Theorem 4]{Tate}, each such class is the class of a
  divisor, so \(M_I\) would contain a nonzero Lefschetz class — contradicting
  that \(\mathcal{M}_I\) is exotic.
\end{proof}

\begin{definition}[Exotic subset]
  \label{def:exotic-subset}
  \lean{MR4199442StandardConjecturesForAbelianFourfolds.ExoticSubset}
  \uses{lem:exotic-frobenius}
  % SOURCE: [Ancona], §7, Definition following (musica) (read from references/source/fourfolds_2020.tex)
  % SOURCE QUOTE: "A subset $I$ of $\Sigma$ of cardinality $4$ verifying (\ref{eq:1}) and (\ref{eq:2}) is called an exotic subset."
  \textit{Source: Ancona, arXiv:1806.03216, §7.}
  A subset \(I \subset \Sigma\) of cardinality four satisfying the two relations
  of \cref{lem:exotic-frobenius}, namely
  \[\prod_{\sigma \in I} \alpha_\sigma = q^2
    \qquad \text{and} \qquad
    \alpha_\sigma \alpha_\tau \neq q
    \ \ (\sigma \neq \tau \in I),\]
  is called an \emph{exotic subset}.
\end{definition}

\begin{remark}
  \label{rem:exotic-completare}
  \uses{lem:exotic-frobenius, def:exotic-subset, rem:exotic-equal-dim}
  % SOURCE: [Ancona], §7, Remark (completare) (read from references/source/fourfolds_2020.tex)
  % SOURCE QUOTE: "Lemma \ref{musica} implies that the dimension of the space of exotic classes is at most the number of exotic subsets. The  Tate conjecture for $A$ predicts that this inequality should actually be an equality."
  \textit{Source: Ancona, arXiv:1806.03216, §7.}
  \cref{lem:exotic-frobenius} shows that any factor \(M_I\) contributing to an
  exotic motive has \(I\) an exotic subset; hence the dimension of the space of
  exotic classes is at most the number of exotic subsets. (The Tate conjecture
  for \(A\) predicts equality.)
\end{remark}

\subsection*{The lemma chain bounding the dimension}

\begin{lemma}[Partition, conjugates, even dimension]
  \label{lem:exotic-partition}
  \lean{MR4199442StandardConjecturesForAbelianFourfolds.exoticSubsetPartition}
  \uses{def:exotic-subset, prop:cm-decomposition}
  % SOURCE: [Ancona], §7, Lemma (lemmaeven) (read from references/source/fourfolds_2020.tex)
  % SOURCE QUOTE: "The following holds.
  % \begin{enumerate}
  % \item\label{partizione} Let $I$ be an exotic subset, then the pair $(I,\bar{I})$ forms a partition of $\Sigma$.
  % \item Any  $\Gal(\overline{\Q}/\Q)$-conjugate of an exotic subset is again exotic.
  % \item Any exotic motive is even dimensional.
  %   \end{enumerate}"
  \textit{Source: Ancona, arXiv:1806.03216, §7.}
  The following hold.
  \begin{enumerate}
    \item If \(I\) is an exotic subset, then the pair \((I, \bar I)\) is a
      partition of \(\Sigma\).
    \item Any \(\Gal(\bar\Q/\Q)\)-conjugate of an exotic subset is again exotic.
    \item Any exotic motive is even-dimensional.
  \end{enumerate}
\end{lemma}

% SOURCE QUOTE PROOF: "Recall that the Weil conjectures imply $\bar{\alpha}=q/\alpha$.
% Then property (\ref{eq:2}) gives the first part. The second part follows from the very definition of exotic subset.
% Altogether complex conjugation acts without fixed points on  the $\Gal(\overline{\Q}/\Q)$-conjugates of $I$, hence the cardinality of the orbit of $I$ under the action of   $\Gal(\overline{\Q}/\Q)$ is even. On the other hand, this cardinality is the dimension of $\mathcal{M}_I$  (see  Proposition \ref{primadecomposizione}), this concludes the last part."
\begin{proof}
  \uses{def:exotic-subset, prop:cm-decomposition}
  The Weil conjectures give \(\bar\alpha = q/\alpha\) for every Frobenius
  eigenvalue \(\alpha\). For (1): if some \(\sigma \in I\) had
  \(\bar\sigma \in I\) as well, then \(\alpha_\sigma \alpha_{\bar\sigma}
  = \alpha_\sigma \cdot q/\alpha_\sigma = q\), contradicting the second relation
  of \cref{def:exotic-subset}; since \(|I| = 4 = |\Sigma|/2\) and \(I\) meets no
  conjugate pair, \((I, \bar I)\) is a partition of \(\Sigma\). Part (2) is
  immediate from the definition of an exotic subset, the two defining relations
  being preserved by the \(\Gal(\bar\Q/\Q)\)-action on the eigenvalues. For (3):
  by (1) complex conjugation sends \(I\) to \(\bar I \neq I\), and the same holds
  for every \(\Gal(\bar\Q/\Q)\)-conjugate of \(I\); thus complex conjugation acts
  without fixed points on the orbit of \(I\), so the orbit has even cardinality.
  By \cref{prop:cm-decomposition} this cardinality equals
  \(\dim \mathcal{M}_I\), so \(\mathcal{M}_I\) is even-dimensional.
\end{proof}

\begin{lemma}[Frobenius relations over extensions]
  \label{lem:exotic-frobenius-power}
  \lean{MR4199442StandardConjecturesForAbelianFourfolds.exoticFrobeniusPower}
  \uses{def:exotic-subset, def:exotic-assumption, rem:exotic-assumption-lef}
  % SOURCE: [Ancona], §7, Lemma after (lemmaeven) (read from references/source/fourfolds_2020.tex)
  % SOURCE QUOTE: "Under the Assumption \ref{assumassum}, an exotic subset $I$ verifies
  % \begin{equation} \label{eq:2n}
  % (\alpha_\sigma \cdot \alpha_\tau)^n \neq q^n, \, \, \forall \sigma \neq \tau, \,\, \sigma,\tau \in I
  % \end{equation}
  % (for all positive integer $n$)."
  \textit{Source: Ancona, arXiv:1806.03216, §7.}
  Under \cref{def:exotic-assumption}, an exotic subset \(I\) satisfies
  \[(\alpha_\sigma \cdot \alpha_\tau)^n \neq q^n
    \qquad \text{for all distinct } \sigma, \tau \in I
    \text{ and all integers } n \ge 1.\]
\end{lemma}

% SOURCE QUOTE PROOF: "If we extend the scalars to  $\mathbb{F}_{q^n}$  the motive $M_I$  will still be exotic (Remark \ref{assumptionlef}).  The eigenvalues of  Frobenius become $\{\alpha_\sigma^n\}$. Hence, by applying  (\ref{eq:2}) over this new base field we deduce (\ref{eq:2n})."
\begin{proof}
  \uses{def:exotic-subset, rem:exotic-assumption-lef}
  Extend scalars to \(\F_{q^n}\). By \cref{rem:exotic-assumption-lef} the motive
  \(M_I\) remains exotic, and the Frobenius eigenvalues become
  \(\{\alpha_\sigma^n\}\). Applying the second relation of
  \cref{def:exotic-subset} over the new base field \(\F_{q^n}\) (whose cardinality
  is \(q^n\)) gives \((\alpha_\sigma \alpha_\tau)^n \neq q^n\) for all distinct
  \(\sigma, \tau \in I\).
\end{proof}

\begin{lemma}[No intersection in exactly two elements]
  \label{lem:exotic-no-two}
  \lean{MR4199442StandardConjecturesForAbelianFourfolds.exoticSubsetsNotTwo}
  \uses{def:exotic-subset, def:exotic-assumption, lem:exotic-frobenius-power}
  % SOURCE: [Ancona], §7, Lemma (lemmanotwo) (read from references/source/fourfolds_2020.tex)
  % SOURCE QUOTE: "Under the Assumption \ref{assumassum}, two exotic subsets cannot intersect in exactly two elements."
  \textit{Source: Ancona, arXiv:1806.03216, §7.}
  Under \cref{def:exotic-assumption}, two exotic subsets cannot intersect in
  exactly two elements.
\end{lemma}

% SOURCE QUOTE PROOF: "Suppose that there are two exotic subsets  $I$ and $J$ whose intersection has cardinality two. Let us call $\alpha,\beta$ the two Frobenius eigenvalues corresponding to  $I-I\cap J$ and $\gamma,\delta$  those corresponding to  $J-I\cap J$.  Property (\ref{eq:1})  gives
%  \[\alpha\cdot \beta=\gamma \cdot  \delta\]
% and  property (\ref{eq:2}) gives (after reordering if necessary)
% \[\alpha= q/ \gamma \,\, , \,\, \beta=  q/\delta.\]
% Putting these relations together we have
% $\alpha^2 =q^2/\beta^2.$ This equality gives a contradiction by applying   (\ref{eq:2n})   to $I$ for $n=2$."
\begin{proof}
  \uses{def:exotic-subset, lem:exotic-frobenius-power}
  Suppose \(I\) and \(J\) are exotic subsets with \(|I \cap J| = 2\). Let
  \(\alpha, \beta\) be the Frobenius eigenvalues attached to \(I \setminus
  (I \cap J)\) and \(\gamma, \delta\) those attached to \(J \setminus (I \cap
  J)\). The first relation of \cref{def:exotic-subset} applied to \(I\) and to
  \(J\) (whose products both equal \(q^2\) and share the factors over
  \(I \cap J\)) gives \(\alpha\beta = \gamma\delta\). The second relation gives,
  after reordering, \(\alpha = q/\gamma\) and \(\beta = q/\delta\). Substituting,
  \(\gamma\delta = (q/\alpha)(q/\beta) = q^2/(\alpha\beta)\), so
  \(\alpha\beta = q^2/(\alpha\beta)\), i.e. \((\alpha\beta)^2 = q^2\). Since
  \(\alpha = \alpha_\sigma\) and \(\beta = \alpha_\tau\) are the eigenvalues of
  the two elements \(\sigma, \tau\) of \(I \setminus (I \cap J)\), this reads
  \((\alpha_\sigma \alpha_\tau)^2 = q^2\), contradicting
  \cref{lem:exotic-frobenius-power} applied to \(I\) with \(n = 2\).
\end{proof}

\begin{lemma}[There are zero, two or four exotic subsets]
  \label{lem:exotic-count}
  \lean{MR4199442StandardConjecturesForAbelianFourfolds.exoticSubsetCount}
  \uses{def:exotic-subset, def:exotic-assumption, lem:exotic-partition,
    lem:exotic-no-two, rem:exotic-completare, rem:exotic-equal-dim}
  % SOURCE: [Ancona], §7, Lemma (lemmaabove) (read from references/source/fourfolds_2020.tex)
  % SOURCE QUOTE: "Under the Assumption \ref{assumassum}, the dimension of the space of exotic classes is zero, two or four. More precisely, there are either  zero, two or four exotic subsets. In case they are two, they are  complex conjugate to each other. In case they are four, they are of the form $I,\bar{I}, J, \bar{J}$, with $I\cap J$ of cardinality three."
  \textit{Source: Ancona, arXiv:1806.03216, §7.}
  Under \cref{def:exotic-assumption}, the dimension of the space of exotic
  classes is zero, two or four. More precisely, there are either zero, two or
  four exotic subsets. If there are two, they are complex conjugate to one
  another. If there are four, they are of the form \(I, \bar I, J, \bar J\) with
  \(|I \cap J| = 3\).
\end{lemma}

% SOURCE QUOTE PROOF: "The statement on exotic classes is implied by the statement on exotic subsets. Indeed, the dimension of the space of exotic classes is even (Lemma \ref{lemmaeven} together with the relation (\ref{eq:uguale}) in Remark \ref{remarkuguale}) and it is  at most the number of exotic subsets (Remark \ref{completare}).
% Let us now show the statement on exotic subsets. We have already pointed out in the proof of Lemma \ref{lemmaeven} that  complex conjugation acts without fixed points on the exotic subsets. Suppose now that there are at least four such subsets, call them $I,\bar{I}, J, \bar{J}$. Then one subset among $J$ and $\bar{J}$ intersects $I$ in at least two elements. Without loss of generality we suppose it is $J$. Then by Lemma \ref{lemmanotwo}, $I\cap J$ must be of cardinality three.
% If there were more than four exotic subsets, then, by the same arguments there would be an exotic subset $K$ intersecting $I$ in exactly three elements. Then the intersection of $J$ and $K$ would have exactly two elements. This is impossible  by Lemma \ref{lemmanotwo}."
\begin{proof}
  \uses{lem:exotic-partition, lem:exotic-no-two, rem:exotic-completare,
    rem:exotic-equal-dim}
  The statement on exotic classes follows from that on exotic subsets: by
  \cref{lem:exotic-partition} together with the equality of
  \cref{rem:exotic-equal-dim} the dimension of the space of exotic classes is
  even, and by \cref{rem:exotic-completare} it is at most the number of exotic
  subsets.

  For the exotic subsets: complex conjugation acts on them without fixed points
  (proof of \cref{lem:exotic-partition}), so they come in conjugate pairs.
  Suppose there are at least four, say \(I, \bar I, J, \bar J\). One of \(J\),
  \(\bar J\) meets \(I\) in at least two elements; say it is \(J\). Since
  \(|I| = |J| = 4\) and \(\bar I = \Sigma \setminus I\), \(J\) cannot equal \(I\)
  or \(\bar I\), so \(1 \le |I \cap J| \le 3\); by \cref{lem:exotic-no-two}
  \(|I \cap J| \neq 2\), and as \(|I \cap J| \ge 2\) we get \(|I \cap J| = 3\).
  If there were more than four exotic subsets, the same argument would produce a
  further exotic subset \(K\) meeting \(I\) in exactly three elements; then
  \(J\) and \(K\) would meet in exactly two elements, impossible by
  \cref{lem:exotic-no-two}. Hence there are zero, two or four exotic subsets,
  with the stated structure.
\end{proof}

\subsection*{The four-dimensional case: reduction to a supersingular product}

\begin{lemma}[Four exotic subsets force a supersingular factor]
  \label{lem:exotic-supersingular}
  \lean{MR4199442StandardConjecturesForAbelianFourfolds.exoticFourSupersingular}
  \uses{def:exotic-subset, def:exotic-assumption, lem:exotic-count,
    lem:exotic-partition, lem:exotic-frobenius}
  % SOURCE: [Ancona], §7, Lemma (lemmafour) (read from references/source/fourfolds_2020.tex)
  % SOURCE QUOTE: "Suppose that there are four exotic subsets and that the Assumption   \ref{assumassum} is satisfied. Then, after extending $k= \mathbb{F}_{q}$ to its quadratic extension $\mathbb{F}_{q^2}$, the abelian fourfold $A$ becomes isogenous to  $E\times X$ where  $X$ is an abelian threefold and $E$ is a supersingular elliptic curve on which   Frobenius acts as $q\cdot \id $."
  \textit{Source: Ancona, arXiv:1806.03216, §7.}
  Suppose there are four exotic subsets and \cref{def:exotic-assumption} holds.
  Then, after extending \(k = \F_q\) to its quadratic extension \(\F_{q^2}\), the
  abelian fourfold \(A\) becomes isogenous to \(E \times X\), where \(X\) is an
  abelian threefold and \(E\) is a supersingular elliptic curve on which
  Frobenius acts as \(q \cdot \id\).
\end{lemma}

% SOURCE QUOTE PROOF: "We keep notation from Lemma \ref{lemmaabove}. If the element of $I-I\cap J$ is $\sigma$ then the element of  $(J-I\cap J)= J \cap \bar{I}$ must be $\bar{\sigma}$ because of  Lemma \ref{lemmaeven}(\ref{partizione}).  Let $\alpha$ be the Frobenius eigenvalue for the action on $M_\sigma$, then  (\ref{eq:1}) applied to $I$ and $J$ implies $\alpha=\bar{\alpha}$, hence $\alpha^2=q$.
% Let us now extend the field of definition to $\mathbb{F}_{q^2}$. Among the eight eigenvalues of  Frobenius we will find $q$, which means that there is a nonzero Frobenius-equivariant map between the Tate module of the supersingular elliptic curve $E$ and the Tate module of $A$. This implies the statement by  \cite[Theorem 4]{Tate}."
\begin{proof}
  \uses{lem:exotic-count, lem:exotic-partition, lem:exotic-frobenius}
  Keep the notation of \cref{lem:exotic-count}: the four exotic subsets are
  \(I, \bar I, J, \bar J\) with \(|I \cap J| = 3\). Let \(\sigma\) be the single
  element of \(I \setminus (I \cap J)\). By \cref{lem:exotic-partition}(1),
  \(\bar J = \Sigma \setminus J\), so \(J \cap \bar I = \bar J \setminus
  \overline{(I \cap J)}\) reduces the element of \(J \setminus (I \cap J)\) to
  \(\bar\sigma\). Writing \(\alpha = \alpha_\sigma\) and using
  \(\bar\alpha = \alpha_{\bar\sigma}\), the first relation of
  \cref{lem:exotic-frobenius} applied to \(I\) and to \(J\) (whose products over
  \(I \cap J\) agree) forces \(\alpha = \bar\alpha\), hence
  \(\alpha^2 = \alpha \bar\alpha = q\).

  Extend the base field to \(\F_{q^2}\); the Frobenius eigenvalues become the
  squares, and \(\alpha^2 = q\) shows that \(q\) occurs among the eight
  eigenvalues. Thus there is a nonzero Frobenius-equivariant map between the
  Tate module of a supersingular elliptic curve \(E\) on which Frobenius acts as
  \(q \cdot \id\) and the Tate module of \(A\). By \cite[Theorem 4]{Tate} this
  yields an isogeny \(A \sim E \times X\) with \(X\) an abelian threefold.
\end{proof}

\begin{lemma}[The threefold factor has dimension two]
  \label{lem:exotic-threefold-factor}
  \lean{MR4199442StandardConjecturesForAbelianFourfolds.exoticThreefoldFactor}
  \uses{def:exotic-subset, def:exotic-assumption, lem:exotic-partition,
    lem:exotic-count, prop:cm-decomposition}
  % SOURCE: [Ancona], §7, Lemma after (lemmafour) (read from references/source/fourfolds_2020.tex)
  % SOURCE QUOTE: "We keep Notation \ref{notnotnot} and Assumption \ref{assumassum}. Consider on the base field $k=\mathbb{F}_{q^2}$   a supersingular elliptic curve $E$  on which   Frobenius acts as $q\cdot \id$.
  % Let $A$ be an abelian fourfold of the form $ A=E\times X$, let  $\Sigma_X$ and $ \Sigma_E$ be CM-structures for $X$ and $E$ and define  $\Sigma=\Sigma_X \cup \Sigma_E$. Then, any exotic subset $I$ verifies that $I\cap \Sigma_X$ has cardinality three. Moreover,  the motive $\mathcal{M}_{I\cap \Sigma_X} \in\CHM (k)_{\Q}$, direct factor of $\mathfrak{h}^3(X)$ as constructed in Proposition \ref{primadecomposizione}, is of dimension two.
  % Finally, if the space of exotic classes on $A$ is four dimensional, it is contained in the four dimensional motive $\mathcal{M}_{I\cap \Sigma_X} \otimes \mathfrak{h}^1(E)$."
  \textit{Source: Ancona, arXiv:1806.03216, §7.}
  Over \(k = \F_{q^2}\), let \(E\) be a supersingular elliptic curve on which
  Frobenius acts as \(q \cdot \id\) and let \(A = E \times X\) with \(X\) an
  abelian threefold. Let \(\Sigma_X\), \(\Sigma_E\) be CM-structures for \(X\),
  \(E\) and put \(\Sigma = \Sigma_X \cup \Sigma_E\). Then any exotic subset \(I\)
  has \(|I \cap \Sigma_X| = 3\), and the motive
  \(\mathcal{M}_{I \cap \Sigma_X} \in \CHM(k)_{\Q}\), a direct factor of
  \(\hh^3(X)\) constructed by \cref{prop:cm-decomposition}, has dimension two.
  Moreover, if the space of exotic classes on \(A\) is four-dimensional, it is
  contained in the four-dimensional motive
  \(\mathcal{M}_{I \cap \Sigma_X} \otimes \hh^1(E)\).
\end{lemma}

% SOURCE QUOTE PROOF: "First recall that  the pair $(I,\bar{I})$ forms a partition of $\Sigma$ by Lemma \ref{lemmaeven}(\ref{partizione}). As $ \Sigma_X$ has cardinality six and is stable by the action of complex conjugation one must have that $I\cap \Sigma_X$ has cardinality three.
% Consider the subsets $K \subset \Sigma_X$ (of cardinality three) verifying
% \begin{equation} \label{eq:4}
% \prod_{\sigma \in K} \alpha_\sigma =q^3
% \end{equation}
% \begin{equation} \label{eq:5}
%  \alpha_\sigma \cdot \alpha_\tau \neq q^2, \, \, \forall \sigma \neq \tau, \,\, \sigma,\tau \in K
% \end{equation}
% and
% \begin{equation} \label{eq:5epoi}
%  \alpha_\sigma  \neq q,  \,\, \forall \sigma\in K
% \end{equation}
% Clearly the relations (\ref{eq:1}) and (\ref{eq:2}) for $I$ imply that $K=I\cap \Sigma_B$ verifies (\ref{eq:4}),(\ref{eq:5}) and (\ref{eq:5epoi}). Conversely, the relations (\ref{eq:4}),(\ref{eq:5}) and (\ref{eq:5epoi})  for $K$ imply that $I=K\cup \{\sigma\}$ verifies (\ref{eq:1}) and (\ref{eq:2}) for any $\sigma \in \Sigma_E$. As there are at most four exotic subsets for $A$, there must be at most two subsets of $\Sigma_X$ verifying  (\ref{eq:4}),(\ref{eq:5}) and (\ref{eq:5epoi}).  On the other hand $I\cap \Sigma_X$  and $\bar{I}\cap \Sigma_X$ are two of those, which implies that there are exactly  two subsets of $\Sigma_X$  verifying (\ref{eq:4}),(\ref{eq:5}) and (\ref{eq:5epoi}).
% Finally, as the Galois group $\Gal(\bar{\Q}/\Q)$ acts on those subsets then the motive $\mathcal{M}_{I\cap \Sigma_X}$ has dimension two.
% The rest follows from the construction of $\mathcal{M}_{I\cap \Sigma_X}$."
\begin{proof}
  \uses{lem:exotic-partition, lem:exotic-count, prop:cm-decomposition}
  By \cref{lem:exotic-partition}(1) the pair \((I, \bar I)\) partitions
  \(\Sigma\). Since \(\Sigma_X\) has cardinality six and is stable under complex
  conjugation, \(I\) meets \(\Sigma_X\) in exactly half of it, so
  \(|I \cap \Sigma_X| = 3\).

  Consider the cardinality-three subsets \(K \subset \Sigma_X\) satisfying
  \[\prod_{\sigma \in K} \alpha_\sigma = q^3, \qquad
    \alpha_\sigma \alpha_\tau \neq q^2 \ (\sigma \neq \tau \in K), \qquad
    \alpha_\sigma \neq q \ (\sigma \in K).\]
  The defining relations of an exotic subset for \(I\) imply that
  \(K = I \cap \Sigma_X\) satisfies these three relations (the factor over
  \(\Sigma_E\) being a single eigenvalue equal to \(q\), since Frobenius acts on
  \(\hh^1(E)\) as \(q\cdot\id\)). Conversely, if \(K \subset \Sigma_X\) satisfies
  them, then \(I = K \cup \{\sigma\}\) is an exotic subset for any
  \(\sigma \in \Sigma_E\). Since \(A\) has at most four exotic subsets
  (\cref{lem:exotic-count}), there are at most two such \(K\); as
  \(I \cap \Sigma_X\) and \(\bar I \cap \Sigma_X\) are two of them, there are
  exactly two. The group \(\Gal(\bar\Q/\Q)\) acts on these two subsets, so the
  motive \(\mathcal{M}_{I \cap \Sigma_X}\) — spanned by the \(M_{g(I \cap
  \Sigma_X)}\) — has dimension two. The final assertion follows from the
  construction of \(\mathcal{M}_{I \cap \Sigma_X}\) in
  \cref{prop:cm-decomposition}.
\end{proof}

\begin{lemma}[Splitting the exotic motive into rank-two pieces]
  \label{lem:exotic-split}
  \lean{MR4199442StandardConjecturesForAbelianFourfolds.exoticSplit}
  \uses{prop:cm-decomposition, cor:lifting-motives, def:hyodo-kato,
    def:abelian-type-rank}
  % SOURCE: [Ancona], §7, Lemma (bellaidea) (read from references/source/fourfolds_2020.tex)
  % SOURCE QUOTE: "Consider an abelian fourfold of the form $ A=E\times X$, where  $E$ is a supersingular elliptic curve such that $\dim_\Q\End(E)\otimes_\Z \Q=4$. Fix a CM-structure for $X$ and consider a direct factor of $\mathfrak{h}^3(X)$  of the form $\mathcal{M}_{I} \in\CHM (k)_{\Q}$, as constructed in Proposition \ref{primadecomposizione}. Suppose that $\mathcal{M}_{I} $ has rank two and that  the  motive ${\mathcal{M}_I \otimes \mathfrak{h}^1(E)}$ is spanned by algebraic classes. Then  there exists a CM structure for  $E$ (write $\Sigma_E=\{\sigma,\bar{\sigma}\}$) such that  the motive $\mathcal{M}_I \otimes \mathfrak{h}^1(E)$ decomposes in the category $\CHM (k)_{\Q}$ as the sum of two motives of rank two
  % \[\mathcal{M}_I \otimes \mathfrak{h}^1(E)= \mathcal{M}_{I\cup \sigma} \oplus \mathcal{M}_{I\cup \bar{\sigma}}\]
  % via Proposition  \ref{primadecomposizione}(\ref{decomposizioneQ})."
  \textit{Source: Ancona, arXiv:1806.03216, §7.}
  Let \(A = E \times X\) with \(E\) a supersingular elliptic curve such that
  \(\dim_\Q \End(E) \otimes_\Z \Q = 4\). Fix a CM-structure for \(X\) and a direct
  factor \(\mathcal{M}_I \in \CHM(k)_{\Q}\) of \(\hh^3(X)\) constructed by
  \cref{prop:cm-decomposition}. Suppose \(\mathcal{M}_I\) has rank two and that
  \(\mathcal{M}_I \otimes \hh^1(E)\) is spanned by algebraic classes. Then there
  is a CM-structure for \(E\), with \(\Sigma_E = \{\sigma, \bar\sigma\}\), such
  that
  \[\mathcal{M}_I \otimes \hh^1(E) =
    \mathcal{M}_{I \cup \sigma} \oplus \mathcal{M}_{I \cup \bar\sigma}\]
  in \(\CHM(k)_{\Q}\), a sum of two motives of rank two, via
  \cref{prop:cm-decomposition}(2).
\end{lemma}

% SOURCE QUOTE PROOF: "By construction of   $\mathcal{M}_I$ (Proposition  \ref{primadecomposizione}(\ref{decomposizioneQ})) we can find a quadratic number field $F$  such that the motive $\mathcal{M}_{I} $  decomposes in the category $\CHM (k)_{F}$ into the sum
% \[\mathcal{M}_{I} = M_{I} \oplus  M_{\bar{I}}  \] of two motives of rank one. In order to conclude, it is enough to show that also $\mathfrak{h}^1(E)$  decomposes in the category $\CHM (k)_{F}$ into the sum of two motives of dimension one.
% This fact is equivalent to saying that $F$ can be embedded in the division algebra $\End(E)\otimes_\Z \Q.$ As this algebra splits at all primes except infinity and  the characteristic of the base field (call it $p$), we have to show that the embeddings $F \subset \R$ and $F\subset \Q_p$ cannot exist. Indeed, as the degree of $F$ over $\Q$ is the same as the index of the central $\Q$-algebra $\End(E)\otimes_\Z \Q$ (namely two), the inclusion  $F\subset\End(E)\otimes_\Z \Q$ is equivalent to the fact that $\End(E)\otimes_\Z F$ is a matrix algebra. By the Albert--Brauer--Hasse--Noether theorem this can be checked locally. As $\End(E)\otimes_\Z \Q_\ell$ is already a matrix algebra it is enough to check at $p$ and at infinity. Now if $F $ does not embed in $\R$ and $\Q_p$ then $F\otimes_{\Q} \R$ and $F\otimes_{\Q}\Q_p$ are quadratic field extensions of $\R$ and $\Q_p$. As $\End(E)\otimes_\Z \Q$ has invariant $1/2$ at $p$ and infinity it becomes a matrix algebra after tensoring with these quadratic extensions.
% We will give two proofs of the non existence of the embeddings $F \subset \R$ and $F\subset \Q_p$.
% A first proof uses that these motives lift to characteristic zero (Proposition \ref{lifting motives}).
% Now if $F$ was contained in $\R$ then the Betti realization of the lifting of $M_{I}$ would respect the Hodge symmetry. This is impossible as it is one dimensional and of weight three. Similarly, if $F$ was contained in $\Q_p$ then the Hyodo-Kato realization of $M_{I}$ and $M_{\bar{I}} $ would be two filtered $\varphi$-modules of dimension one with different filtrations (again because the weight of $\mathcal{M}_{I}$ is three). This implies that
% $\det (M_{I}  \otimes \mathfrak{h}^1(E))$ and $\det(M_{\bar{I}}  \otimes \mathfrak{h}^1(E))$ would realize in two filtered $\varphi$-modules of dimension one and different filtrations.  On the other hand   absolute Frobenius acts on both   in the same way (namely by multiplication by $p^4$) because they are spanned by algebraic classes. This implies that at least one of the two filtered $\varphi$-modules is not admissible and concludes the proof as all filtered $\varphi$-modules coming from geometry (in particular filtered $\varphi$-modules that are realizations of motives) must be admissible.
% We give an alternative proof, which does not use the lifting to characteristic zero. We write it for $\Q_p$, it works in the same way for $\R$. If $F$ was contained in $\Q_p$ then the motive $M_{I}  \otimes \mathfrak{h}^1(E)$  would live in $\CHM (k)_{\Q_p}$. On such a motive the division algebra $\End(E)\otimes_\Z \Q_p$ acts hence it acts also on the $\Q_p$ vector space spanned by the algebraic classes of the motive. On the other hand, as this motive is spanned by algebraic classes, the division algebra  $\End(E)\otimes_\Z \Q_p$  would act on a $\Q_p$-vector space  of dimension two. This gives a contradiction as such an  action cannot exist."
\begin{proof}
  \uses{prop:cm-decomposition, cor:lifting-motives, def:hyodo-kato}
  By the construction of \(\mathcal{M}_I\) (\cref{prop:cm-decomposition}(2)) there
  is a quadratic number field \(F\) over which \(\mathcal{M}_I\) splits in
  \(\CHM(k)_{F}\) as a sum \(\mathcal{M}_I = M_I \oplus M_{\bar I}\) of two
  rank-one motives. It suffices to show that \(\hh^1(E)\) likewise splits in
  \(\CHM(k)_{F}\) into a sum of two motives of dimension one; this is equivalent
  to embedding \(F\) into the division algebra \(\End(E) \otimes_\Z \Q\). The
  quaternion algebra \(\End(E) \otimes_\Z \Q\) is split at every prime except
  \(p\) and \(\infty\) (where its invariant is \(1/2\)); by the
  Albert--Brauer--Hasse--Noether theorem the embedding exists iff
  \(\End(E) \otimes_\Z F\) is a matrix algebra, which need only be checked at
  \(p\) and \(\infty\). Thus it suffices to rule out the embeddings
  \(F \subset \R\) and \(F \subset \Q_p\).

  Both are excluded using that the motives lift to characteristic zero
  (\cref{cor:lifting-motives}). If \(F \subset \R\), the Betti realization of the
  lift of \(M_I\) would have to respect Hodge symmetry, impossible for a
  one-dimensional space of weight three. If \(F \subset \Q_p\), the Hyodo--Kato
  realizations (\cref{def:hyodo-kato}) of \(M_I\) and \(M_{\bar I}\) would be
  one-dimensional filtered \(\varphi\)-modules with distinct filtrations (the
  weight being three), whence \(\det(M_I \otimes \hh^1(E))\) and
  \(\det(M_{\bar I} \otimes \hh^1(E))\) would be one-dimensional filtered
  \(\varphi\)-modules with distinct filtrations on which absolute Frobenius acts
  identically (by \(p^4\), as both are spanned by algebraic classes); one of them
  would then fail to be admissible, contradicting that realizations of motives
  are admissible. (Alternatively, without lifting: if \(F \subset \Q_p\) the
  division algebra \(\End(E) \otimes_\Z \Q_p\) would act on the two-dimensional
  \(\Q_p\)-space of algebraic classes of \(M_I \otimes \hh^1(E)\), which is
  impossible.) Hence \(F\) embeds in \(\End(E) \otimes_\Z \Q\), \(\hh^1(E)\)
  splits over \(F\), and the two rank-two summands
  \(\mathcal{M}_{I \cup \sigma}, \mathcal{M}_{I \cup \bar\sigma}\) are obtained by
  \cref{prop:cm-decomposition}(2).
\end{proof}

\subsection*{Proof of the main proposition}

The proof of \cref{prop:exotic-dim2} was given immediately after its statement at
the start of this chapter; it combines the lemmas
\cref{lem:exotic-count}, \cref{lem:exotic-partition},
\cref{lem:exotic-supersingular}, \cref{lem:exotic-threefold-factor} and
\cref{lem:exotic-split} established above.

\section{Geometric examples}

\subsection*{No exotic lift of Betti type \texorpdfstring{$(4,0),(0,4)$}{(4,0),(0,4)}}

\begin{remark}[No exotic lift of Betti type \((4,0),(0,4)\)]
  \label{rem:exotic-no-top-type}
  \uses{def:exotic-motive, def:lefschetz-class, cor:lifting-motives,
    prop:decomposition, lem:exotic-supersingular}
  % SOURCE: [Ancona], Appendix A, Remark (altreesotiche)(altreesotiche2) (read from references/source/fourfolds_2020.tex)
  % SOURCE QUOTE: "There are no exotic motives over $\overline{\mathbb{F}}_p$ (coming from abelian fourfolds) whose lifting  to $\C$ have Betti realization    of type $(4,0),(0,4)$. To show this, consider a model of the abelian fourfold over a finite field $\mathbb{F}_{q}$. Let $I$ be the set of Frobenius eigenvalues such that the corresponding eigenspaces are lifted into $H^{1,0}$ and $\bar{I}$ be the set of Frobenius eigenvalues such that the corresponding eigenspaces are lifted into $H^{0,1}$. If the cohomology group $H^{4,0}\oplus H^{0,4}$ becomes Frobenius invariant over $\mathbb{F}_q$, then $\prod_{\alpha\in I} \alpha = \prod_{\beta\in \bar{I}} \beta(=q^2).$ On the other hand, using the Shimura--Taniyama formula  \cite[Lemma 5]{TateBour}, we have that the $p$-adic valuation of  $\prod_{\alpha\in I} \alpha$ is greater than the one of $ \prod_{\beta\in \bar{I}} \beta$, except if all Frobenius eigenvalues have the same slopes. In this case the abelian variety would be isogenous to the forth power of a supersingular elliptic curve and hence all algebraic classes would be Lefschetz."
  \textit{Source: Ancona, arXiv:1806.03216, Appendix A.}
  There are no exotic motives over \(\overline{\F}_p\) (coming from abelian
  fourfolds) whose lifting to \(\C\) has Betti realization of type
  \((4,0),(0,4)\). In particular, for the application to abelian fourfolds only
  the case \(i = 1\) occurs.

  To see this, take a model of the abelian fourfold over a finite field
  \(\F_{q}\). Let \(I\) be the set of Frobenius eigenvalues whose eigenspaces lift
  into \(H^{1,0}\) and \(\bar I\) the set of those whose eigenspaces lift into
  \(H^{0,1}\). If \(H^{4,0} \oplus H^{0,4}\) becomes Frobenius invariant over
  \(\F_q\), then
  \[\prod_{\alpha \in I} \alpha = \prod_{\beta \in \bar I} \beta \, (= q^2).\]
  On the other hand, by the Shimura--Taniyama formula \cite[Lemma 5]{TateBour},
  the \(p\)-adic valuation of \(\prod_{\alpha \in I} \alpha\) is strictly greater
  than that of \(\prod_{\beta \in \bar I} \beta\), unless all Frobenius
  eigenvalues have the same slope. In that exceptional case the abelian variety
  would be isogenous to the fourth power of a supersingular elliptic curve, so all
  algebraic classes would be Lefschetz; hence no exotic motive of this type can
  occur.
\end{remark}
